\section{模型假设}\label{sec:model-assumptions}

题目附录已经给出了设备参数和动作规则，本文直接沿用，不再重复。除此之外，建模时还用到下面几条简化和约定。前六条对四个问题都适用，后两条只在第四问用到。

\begin{enumerate}
  \item \textbf{平面运动。}忽略机器狗与干扰源的高度差，把两者都看作平面上的点。机器狗知道自己的准确位置，能走到任何指定的坐标；两点之间走直线，速度恒为 $5$ m/s，不考虑加减速、转向和绕行，因此移动时间等于欧氏距离除以速度。

  \item \textbf{干扰源静止，相互独立。}一次任务中干扰源不移动，被清除之前它的频道、发射功率和接收半径都不变。不同源的信号互不影响，清除一个源不改变其他源能否被收到。源的真实位置、个数和接收半径算法事先都不知道，只能通过检测和清除的反馈逐步了解。

  \item \textbf{能否收到信号只由距离决定。}在源的发射范围之内，检测点与源的距离不超过接收半径就一定能收到信号，超过就一定收不到，不考虑遮挡、间歇发射和随机漏检。距离不超过 $5$ m 时按题设返回近场反馈，不再给出示向度。全向源的发射范围是整个平面，所以在第二、三问中，某处收不到信号就说明该处离源超过了接收半径。

  \item \textbf{测角误差只当作有界的不确定量。}示向度误差不超过 $\pm1^\circ$，除此之外不对误差的分布作任何假定，也不假定不同地点的误差相互独立；同一地点重复检测不能缩小误差。程序里把误差界取为 $1.005^\circ$，多出的 $0.005^\circ$ 用来容纳示向度保留两位小数带来的舍入，定位区域仍然只按真实反馈收缩。

  \item \textbf{动作顺序执行，按题设计费。}移动、换频、检测、光学定位和清除依次进行，不能同时做，任务总耗时是各项动作耗时之和。光学定位失败同样计 $3$ s，成功后再计 $2$ s 的清除时间。程序的计算时间另行记录，不计入机器狗的动作耗时。

  \item \textbf{光学定位与清除只看距离。}机器狗到源的距离不超过 $20$ m 时，光学定位和清除一定成功，成功与否和源的发射方向无关；超过 $20$ m 则一定失败。算法只在收到设备返回的清除成功之后才把该频道记为已清除，没有得到确认的请求不算成功。

  \item \textbf{定向源向半平面辐射。}第四问中的定向源只向发射方向两侧各 $90^\circ$ 的范围辐射，即过源位置、垂直于发射方向的直线所限定的闭半平面；只有位于这个半平面内且距离不超过接收半径的位置才能收到信号，近场反馈同样受此限制。发射方向在一次任务中不变，定向源的个数和所占比例都未知。

  \item \textbf{机器狗可以走出目标区域。}题目只规定干扰源在半径 $1800$ m 的圆内，没有限制机器狗的活动范围。本文允许机器狗到圆外不远处停下检测，用来发现贴着边界朝外发射的定向源；第四问的检测站中有一环放在半径 $1865$ m 的圆周上。
\end{enumerate}

